\subsubsection{使用更少数据进行微调}

指令微调常依赖大规模数据集，例如 FLAN 汇集了 1836 个任务、约 1500 万个样本 \cite{longpre-etal:2023flan}。在如此庞大的数据上对大型语言模型进行全参数更新，计算开销极高。参数高效方法通过只更新少量参数来缓解这一问题，但微调数据本身仍可能包含合成噪声和偏差。另一种高效思路是从数据入手：\textbf{只选用最相关、最有影响力的样本进行微调}，在保持模型更新质量的同时大幅压缩数据量。形式化地，给定大规模数据集 $\mathcal{D}$，我们希望选出一个高信息密度的子集 $\mathcal{S}^* \subset \mathcal{D}$，使基于该子集微调得到的模型 $\theta^*(\mathcal{S}^*)$ 在目标任务上的损失尽可能逼近使用全量数据的效果：
\begin{equation}
\mathcal{S}^* = \arg\min_{\mathcal{S}\subseteq\mathcal{D},\,|\mathcal{S}|=k} \mathcal{L}_{\text{task}}\bigl(\theta^*(\mathcal{S})\bigr).
\end{equation}

多项工作验证了这一思路。\citet{zhou-etal:2023lima} 精心设计提示并收集样本，构建了仅包含 1000 个样本的指令遵循数据集 LIMA。使用该数据微调的 LLaMA 65B 模型，其竞争力可与投入更多微调努力的模型相媲美，甚至更优：在人类偏好评估中，LIMA-65B 对 DaVinci003 的胜率达 58\%，而用 52k 条数据微调的 Alpaca-65B 仅获 47\% 的胜率。这说明模型的多任务响应能力无需在所有指令类型上微调即可激发。\citet{chen-etal:2024alpagasus} 则利用 GPT-3.5 为每条指令-回答对打分，从 Alpaca-52K 中筛选出 9000 个最高质量样本构成 AlpaGasus-9K；仅用该子集微调，在 AlpacaEval 上的表现即持平甚至超过全量数据。这些结果进一步印证了“数据质量优先于数量”的观点。此外，LESS 等方法通过计算每个训练样本对验证损失的梯度影响来筛选最有价值的样本：
\begin{equation}
\mathcal{I}(z) \approx -\nabla_\theta \ell(z)^\top \mathbf{H}^{-1} \nabla_\theta \ell(\mathcal{D}_{\text{val}}),
\end{equation}
其中 $\mathbf{H}$ 为海森矩阵。选择影响力最大的子集，可高效驱动模型收敛。总体而言，使用更少但更高质量的数据训练 NLP 模型，往往能带来更好效果。

上述发现与 NLP 的传统认知有所不同：复杂任务的执行能力可以通过少量标注数据激活，而非依赖大规模监督训练。一种解释是，生成正确响应的能力已在预训练阶段习得，但指令-响应的映射在推理时并非高概率路径。微调只需对模型进行微小调整即可使其遵循指令，所需的训练量远少于预训练。

\subsubsection{参数高效微调方法}

即使有监督微调的数据量远小于预训练，在数十亿参数的模型上进行全参数微调仍消耗可观的 GPU 显存与计算时间。这促使了参数高效微调（Parameter-Efficient Fine-Tuning, PEFT）方法的产生。与更新全部参数的全参数微调（Full Fine-Tuning, FFT）不同，PEFT 冻结预训练模型的大部分权重，仅训练少量附加参数或特定层。早期的 PEFT 方法通过在 Transformer 层之间插入小型全连接网络（\mindex{adapter}）\cite{houlsby-etal:2019adapter}，或在输入端添加可学习的连续提示向量（如 \mindex{prefix tuning} 和 \mindex{prompt tuning}）来引导生成 \cite{li-liang:2021prefix,lester-etal:2021prompt}。但这些方法往往在推理时引入额外延迟，或在复杂任务上难以追平全参数微调的性能上限。

为克服上述局限，\citet{hu-etal:2021lora} 提出了低秩自适应（Low-Rank Adaptation, \mindex{LoRA}）。其核心假设是：微调过程中的权重更新 $\Delta W$ 具有较低的“内在秩”。如图~\ref{fig:lora} 所示，对于预训练权重 $W_0 \in \mathbb{R}^{d \times k}$，LoRA 将其更新分解为两个低秩矩阵的乘积：
\begin{equation}
W' = W_0 + \Delta W = W_0 + B A,\quad \text{其中 } B \in \mathbb{R}^{d \times r},\; A \in \mathbb{R}^{r \times k},\; r \ll \min(d,k).
\end{equation}
训练时 $W_0$ 被冻结，仅优化 $B$ 和 $A$。秩 $r$ 通常取 4、8 或 16，可带来巨大的参数缩减。以 LLaMA-7B 的某个注意力投影矩阵 ($d=k=4096$) 为例，若 $r=16$，原始矩阵参数量约 $4096^2 \approx 16.7\text{M}$，而 LoRA 的参数量仅 $2 \times 4096 \times 16 \approx 131\text{K}$，减少超过 100 倍。

\begin{figure}[htbp]
  \centering
  \begin{tikzpicture}[
    >=Latex,
    every node/.style={font=},
    frozen/.style={rounded corners=6pt, draw=gray!60, fill=gray!10, line width=1pt},
    train/.style={rounded corners=6pt, draw=orange!70!black, fill=orange!12, line width=1pt},
    plusnode/.style={circle, draw=gray!60, fill=gray!10, line width=1pt, minimum size=0.9cm, font=\bfseries},
    iobox/.style={rounded corners=4pt, draw=gray!55, fill=gray!6, line width=0.8pt, minimum width=1.2cm, minimum height=0.9cm},
    sub/.style={font=\scriptsize, gray!40!black},
    arrow/.style={->, gray!45!black, line width=0.8pt},
  ]
  
  % ---- 输入 x ----
  \node[iobox] (x) at (0,0) {$x$};
  \coordinate (split) at (1.3,0);
  \draw[arrow] (x.east) -- (split);
  
  % ---- 上方路径：W0（冻结）----
  \node[frozen, minimum width=3.4cm, minimum height=1.8cm] (w0) at (4.2,1.3) {};
  \node[font=\bfseries] at (4.2,1.52) {$W_0$};
  \node[sub] at (4.2,1.24) {冻结, $d \times k$};
  \node[sub] at (4.2,0.99) {不参与梯度更新};
  
  % ---- 下方路径：先 A（r x k，压缩到秩 r）----
  \node[train, minimum width=1.5cm, minimum height=0.9cm] (a) at (3.6,-1.2) {};
  \node[font=\bfseries] at (3.6,-1.05) {$A$};
  \node[sub] at (3.6,-1.4) {$r \times k$};
  
  % ---- 再 B（d x r，展开回维度 d）----
  \node[train, minimum width=1.5cm, minimum height=1.8cm] (b) at (5.6,-1.2) {};
  \node[font=\bfseries] at (5.6,-0.9) {$B$};
  \node[sub] at (5.6,-1.8) {$d \times r$};
  
  % ---- 分支：x 进入上下两条路径 ----
  \draw[arrow] (split) |- (w0.west);
  \draw[arrow] (split) |- (a.west);
  
  % ---- 下方路径内部：A 先算，结果送入 B ----
  \draw[arrow] (a.east) -- (b.west);
  
  % ---- 汇合点（+）----
  \node[plusnode] (plus) at (8.0,0) {$+$};
  \draw[arrow] (w0.east) -| (plus.north);
  \draw[arrow] (b.east)  -| (plus.south);
  
  % ---- 输出 h ----
  \node[iobox] (h) at (9.6,0) {$h$};
  \draw[arrow] (plus.east) -- (h.west);
  
  \end{tikzpicture}
  \caption{LoRA 示意图。输入 $x$ 从左侧进入，分两条路径向右传播：上方路径穿过冻结的预训练权重 $W_0$；下方路径依次穿过可训练的低秩矩阵 $A$（$r \times k$，先将输入压缩到秩 $r$）和 $B$（$d \times r$，再展开回维度 $d$），两条路径最终在右侧汇合相加，得到输出 $h = W_0x + BAx$。}
  \label{fig:lora}
  \end{figure}

LoRA 的一大优势是训练后低秩矩阵可直接合并回原权重 ($W' = W_0 + BA$)，推理时不会引入任何额外延迟。近年来涌现出众多 LoRA 变体：QLoRA 将 $W_0$ 以 4-bit NormalFloat 量化存储并配合分页优化器，使得在单张 48GB 消费级 GPU 上微调 65B 模型成为可能 \cite{dettmers-etal:2023qlora}；DoRA 则将权重矩阵解耦为幅度与方向分别学习，进一步提升了微调的稳定性和任务表现 \cite{liu-etal:2024dora}。

表~\ref{tab:peft-memory} 对比了全参数微调与主流 PEFT 方法在 7B 模型上的显存占用（混合精度训练）。PEFT 方法将可训练参数量降低数个数量级，显著缓解了显存压力，使大模型微调走向“平民化”。

\begin{table}[htbp]
\centering
\caption{全参数微调与 PEFT 方法的资源对比（以 7B 模型，序列长度 512，批次大小 1 为例）}
\label{tab:peft-memory}
\begin{tabular}{lccc}
\hline
\textbf{方法} & \textbf{可训练参数占比} & \textbf{GPU 显存占用 (估)} & \textbf{推理延迟} \\
\hline
全参数微调 (FFT) & 100\% & $\sim$55 GB & 无额外 \\
Adapter           & 3\%--5\%  & $\sim$20 GB & 略有增加 \\
Prefix/Prompt Tuning & $<$1\% & $\sim$16 GB & 略有增加 \\
LoRA ($r$=16)     & 0.1\%--0.5\% & $\sim$18 GB & \textbf{无增加} \\
QLoRA (4-bit)     & 0.1\%--0.5\% & $\sim$10 GB & 无增加 \\
\hline
\end{tabular}
\end{table}

\begin{importantnote}
LoRA 通过低秩分解将权重更新压缩为两个小矩阵的乘积，仅训练极少量参数即可达到接近全量微调的性能。其核心优势在于推理时零延迟（可合并回原权重），且显存占用大幅降低。需注意秩r的选择会影响表现，通常 8～16 即可，过大则失去高效性。
\end{importantnote}


